% ID: SEL-D0037-OB-E03
% ESTADO: cotejada. Segunda pasada visual y borrador compilado revisados.
% Fuente: Selectividad/matematicas_aplicadas/Ua2MATEMÁTICAS APLICADAS examen.pdf
% SHA-256: a56739d297d4dc220d6898b24cf729dc556401a6bb292b8ca7cfb58c01c31651
% Página: 2; opción B; ejercicio original 3.

Sean dos sucesos $A$ y $B$ tales que $P(A)=0.25$, $P(B)=0.6$, $P(A\cap B^C)=0.1$.
\begin{enumerate}[label=\alph*)]
\item \textbf{(0.75 puntos)} Calcule la probabilidad de que ocurra $A$ y ocurra $B$.
\item \textbf{(0.75 puntos)} Calcule la probabilidad de que no ocurra $A$ pero sí ocurra $B$.
\item \textbf{(0.5 puntos)} Calcule la probabilidad de que ocurra $A$ sabiendo que ha ocurrido $B$.
\item \textbf{(0.5 puntos)} \textquestiondown{}Son independientes $A$ y $B$?
\end{enumerate}
